\section{Introduction}
\label{sec:intro}

Dense pixel-level annotation of 3D medical volumes is one of the most expensive supervision signals in computer vision. A single pancreatic CT volume takes a radiologist tens of minutes to contour; a left-atrium MR volume requires manual correction across hundreds of slices; inter-rater variability on soft-tissue boundaries can exceed the accuracy gap between competing segmentation methods. The practical consequence is that, across every major 3D medical segmentation benchmark, the number of labeled volumes is measured in tens while the pool of unlabeled scans available from the same imaging protocol is measured in hundreds or thousands. Semi-supervised learning (SSL) is therefore not a convenience but the default operating regime for research and production medical segmentation: the relevant question is not whether to exploit unlabeled data but how.

\paragraph{The SSL plateau.}
The dominant SSL methods for 3D medical segmentation --- mean-teacher consistency~\cite{tarvainen2017meanteacher}, uncertainty-aware mean teacher (UA-MT)~\cite{yu2019uamt}, dual-task consistency (DTC)~\cite{luo2022dtc}, bidirectional copy-paste (BCP)~\cite{bai2023bcp}, multi-formation soft masking (MOST)~\cite{liu2024most}, and dynamic uncertainty-aware contrastive learning (DyCON)~\cite{assefa2025dycon} --- differ substantially in how they couple labeled and unlabeled supervision, but they share an implicit contract with the unlabeled pool: training runs for a fixed number of epochs against a composite loss, and at convergence the unlabeled data is discarded. The final released checkpoint is then treated as terminal. This assumption is rarely examined in print, but it is what keeps reported performance on canonical Pancreas-CT and Left Atrium benchmarks stuck within a narrow band below the supervised-with-100\%-labels ceiling, despite architectural improvements, richer augmentations, and increasingly elaborate teacher--student couplings.

\paragraph{Residual uncertainty is not residual noise.}
We revisit the discarded checkpoints with a single question: what does the converged SSL network still \emph{not know}, and is what it doesn't know a usable signal? To answer this, we measure the per-voxel softmax distribution of the publicly released BCP checkpoints on the unlabeled training volumes and decompose the total prediction entropy $\sum_i H_i$ into its contribution from low-confidence voxels (``the uncertain shell'') and from high-confidence voxels (``the saturated interior''). The result is surprisingly sharp and consistent across datasets and label fractions: \textbf{$89$--$95\%$ of the network's total prediction entropy is carried by an uncertain-voxel shell occupying only $5$--$9\%$ of the volume}, while a supervised-only baseline trained on the same $12$ labeled Pancreas-CT volumes carries only $\sim\!55\%$ of its entropy in the corresponding shell (Sec.~\ref{sec:method}, Table~\ref{tab:entropy}). The gap between these two numbers --- $34$ to $40$ percentage points --- is larger than any individual method's reported improvement over the previous year's state of the art on these benchmarks, and it is observed consistently across two datasets and three label fractions with no tuning.

\paragraph{What the measurement means.}
A network whose residual entropy lives at physically meaningful locations (the organ surface) and a network whose residual entropy is scattered across semantically flat interior voxels are in qualitatively different regimes. Sharpening the first is an act of boundary refinement on a near-correct contour: each gradient step slides the predicted surface along a direction in which the labeled examples already established a margin. Sharpening the second is an act of \emph{committing} to guesses on voxels the network does not understand, which is well known to reinforce rather than correct errors~\cite{grandvalet2004entropy,lee2013pseudolabel}. SSL convergence produces the first regime; supervised-only training on a small labeled set produces the second. The concentration of entropy in a boundary shell is, in this sense, the useful residue of SSL training: the signal that tells a post-hoc refinement step where, and only where, there is work to be done.

\paragraph{Post-hoc Entropy Minimization (PEM).}
Motivated by this observation, we introduce \textbf{Post-hoc Entropy Minimization}: given any converged SSL checkpoint, fine-tune it for a small fixed budget ($2$--$5$ epochs) on the unlabeled training set using a single entropy-minimization loss, optionally restricted to the low-confidence shell via an explicit threshold mask. PEM uses no labels, generates no pseudo-labels, introduces no architectural changes, and adds no hyperparameter to the base SSL recipe beyond a dataset-specific learning rate. The entire fine-tuning run completes in under five minutes per configuration on a single GPU. Applied to the publicly released BCP checkpoints without any test-set hyperparameter tuning, PEM improves Dice by $+1.12$, $+1.43$, and $+0.64$ on Pancreas-CT 20\%, LA 5\%, and LA 10\% respectively, and reduces the 95\% Hausdorff distance by $27$--$43\%$. All three Dice gains pass a per-case Wilcoxon signed-rank test at $p<0.025$ with bootstrap $95\%$ confidence-interval lower bound $\geq\!+0.18$ Dice, on test sets of $N\!=\!18$ to $20$ cases.

\paragraph{What is new and what is not.}
The entropy-minimization loss itself is classical and is not our contribution. Grandvalet and Bengio~\cite{grandvalet2004entropy} introduced it as an SSL regularizer over two decades ago; FixMatch~\cite{sohn2020fixmatch} uses a thresholded form of it as part of its consistency objective; entropy-based masking is a standard ingredient in modern SSL. All of these are \emph{training-time} interventions, where the loss acts on a network whose per-voxel distributions are diffuse, still evolving, and coupled to a supervised gradient that prevents collapse. Our contribution is a regime shift: the same loss, applied to a converged optimum on the other side of a long SSL training run, is a qualitatively different object. Because the residual entropy of the converged network is already concentrated at a thin, physically meaningful shell (Table~\ref{tab:entropy}), the gradient signal of the loss is automatically aligned with the locus of error. There is no need for a class-balance regularizer, a warmup curriculum, or a confidence threshold tied to a supervised loss; the $p(1{-}p)$ factor of the binary-entropy gradient, together with the convergence-induced sharpness of the posterior, does the gating implicitly (Sec.~\ref{sec:theory}). The resulting method is a label-free, architecture-free, schedule-free fine-tune that can be applied to any released checkpoint and that inherits decades of literature on the loss itself while occupying a different operating regime than any of its predecessors.

\paragraph{A diagnostic that predicts applicability.}
The same measurement that motivated the method serves as its suitability check. We define the \textbf{shell entropy concentration} $\rho_f(\tau)$ as the fraction of total prediction entropy carried by voxels of confidence at most $\tau$. This quantity requires only a forward pass over unlabeled data, is scalar, and is comparable across networks, datasets, and label fractions. Empirically, the threshold $\rho_f(0.99)\!\ge\!0.85$ separates the regime in which PEM is well-conditioned from the regime in which it is not: on the three BCP checkpoints we study, $\rho_f$ lies in $[0.894, 0.948]$ and PEM improves Dice; on the supervised-only baseline, $\rho_f\!=\!0.552$ and PEM is halted by a collapse-protection safeguard within a single epoch. We recommend reporting $\rho_f(0.99)$ alongside any released SSL checkpoint as a cheap, label-free indicator of whether a refinement step of the kind we propose is expected to help.

\paragraph{Contributions.}
\begin{enumerate}
    \setlength\itemsep{1pt}
    \item \textbf{An empirical regularity of SSL convergence.} The converged BCP checkpoints we study concentrate $89$--$95\%$ of their residual prediction entropy in a thin uncertain-voxel shell occupying $5$--$9\%$ of each volume, whereas a supervised-only baseline trained on the same labels concentrates only $\sim\!55\%$. To the best of our knowledge, this is the first quantitative characterization of the spatial distribution of residual uncertainty in converged SSL segmentation checkpoints on these benchmarks. The regularity holds across two datasets and three label fractions without tuning.
    \item \textbf{A scalar, label-free diagnostic $\rho_f(\tau)$} that quantifies shell entropy concentration and separates the regime in which a post-hoc refinement is well-conditioned from the regime in which it is not. Its computation requires one forward pass per unlabeled volume; we show that $\rho_f(0.99)\!\ge\!0.85$ is a reliable applicability check on the checkpoints we test, and that the ordering of PEM gains across configurations tracks $\rho_f$ and the residual entropy budget in a manner consistent with the theoretical argument of Section~\ref{sec:theory}.
    \item \textbf{Post-hoc Entropy Minimization (PEM).} A label-free, pseudo-label-free, architecture-free fine-tuning step that applies to any converged SSL checkpoint. PEM differs from the classical entropy-minimization literature in its operating regime rather than in its loss: it acts on a converged optimum whose residual entropy is already structurally concentrated, so the gradient is aligned with the locus of error by construction and requires no auxiliary regularization to avoid collapse.
    \item \textbf{A theoretical account} of why this post-hoc regime works while training-time entropy minimization on diffuse distributions does not. The $\partial H/\partial z = -z\,p(1{-}p)$ gradient of the binary entropy is an implicit boundary-shell gate that is only useful when the posterior is already sharp; the $p(1{-}p)$ factor simultaneously prevents collapse, matches the cluster assumption of Chapelle et al.~\cite{chapelle2006ssl} at the physical organ surface, and predicts both failure modes we observe (Sec.~\ref{sec:theory}).
    \item \textbf{Reproducible empirical validation.} Under a fixed-budget protocol with no test-set hyperparameter tuning and three random seeds per configuration, PEM improves Dice by $+0.6$ to $+1.4$ and HD95 by $27$--$43\%$ on Pancreas-CT 20\%, LA 5\%, and LA 10\%, in under $5.2$ minutes per run on a single consumer GPU. A paired per-case analysis (Wilcoxon signed-rank, bias-corrected bootstrap) rejects the null of zero median improvement at $p\!<\!0.025$ on every configuration, with CI lower bound $\geq\!+0.18$ Dice and paired effect sizes $d_z\!\in\![0.57, 0.71]$.
    \item \textbf{Isolation of the mechanism via three controlled post-hoc baselines.} Applied to the same BCP checkpoint under the same fixed-budget protocol, \emph{temperature scaling} (pure recalibration) produces $|\Delta\mathrm{Dice}|\!\le\!0.03$ on every setting, confirming that PEM's gains cannot be attributed to confidence rescaling. \emph{Pseudo-label fine-tuning} (hard argmax targets with connected-component filtering) and \emph{self-consistency} (post-hoc mean-teacher) both produce smaller gains than PEM on the two less-saturated configurations and match PEM in the saturation limit, a tradeoff that is quantitatively predicted by the theoretical argument of Section~\ref{sec:theory}.
    \item \textbf{Characterization of failure modes.} Outside the regime delineated by $\rho_f(0.99)$, PEM fails in two distinct and theoretically predicted ways: \emph{diffuse-entropy collapse} when $\rho_f\!\lesssim\!0.6$ (observed on the supervised baseline) and \emph{vanishing-entropy diminishing returns} as $\rho_f\!\to\!1$ (observed on LA 10\%). Both are ruled out in advance by a single label-free measurement on the unlabeled training set.
\end{enumerate}

The remainder of the paper proceeds as follows. Section~\ref{sec:related} positions PEM against the classical entropy-minimization literature, the SSL medical-segmentation literature, test-time adaptation, calibration, and post-hoc model-refinement methods. Section~\ref{sec:method} defines $\rho_f(\tau)$ and the PEM loss, and reports the entropy-concentration measurements that motivate the method. Section~\ref{sec:theory} provides the theoretical account of why the loss is well-conditioned on a converged checkpoint and ill-conditioned elsewhere, including a derivation of the $p(1{-}p)$ implicit-gating argument and a connection to the cluster assumption at the physical organ surface. Section~\ref{sec:experiments} presents the main results on Pancreas-CT and LA, including the three controlled post-hoc baselines and the per-case statistical analysis. Section~\ref{sec:ablations} reports the ablations over masking threshold, learning rate, and epoch budget. Section~\ref{sec:discussion} discusses failure modes, limitations, and directions for future work.
